\documentclass{article}
\usepackage[utf8]{inputenc}
\usepackage[spanish]{babel}
\usepackage{tikz}
\usepackage{geometry}

\geometry{a4paper, margin=0.5in, landscape}
\usetikzlibrary{shapes.geometric, arrows.meta, positioning, calc}

\begin{document}

\begin{figure}[h]
\centering

\begin{tikzpicture}[
    scale=0.9,
    transform shape,
    % --- ESTILOS (Blanco y Negro Limpio) ---
    entity/.style={
        rectangle,
        draw=black,
        thick,
        fill=white,
        minimum width=2.5cm,
        minimum height=1.2cm,
        align=center,
        font=\bfseries
    },
    relationship/.style={
        diamond,
        aspect=1.5,
        draw=black,
        thick,
        fill=white,
        minimum width=2cm,
        align=center,
        font=\small\itshape,
        inner sep=1pt
    },
    attribute/.style={
        ellipse,
        draw=black,
        thick,
        fill=white,
        minimum width=1.5cm,
        minimum height=0.7cm,
        align=center,
        font=\scriptsize,
        inner sep=1pt
    },
    % clave primaria: mismo estilo que attribute pero SIN \underline
    key_attribute/.style={
        attribute,
        font=\scriptsize\bfseries
    },
    link/.style={-, thick, draw=black},
    card/.style={
        font=\scriptsize\bfseries,
        text=black,
        fill=white,
        inner sep=2pt
    }
]

% ==========================================
% 1. ESTRUCTURA PRINCIPAL (ENTIDADES)
% ==========================================

\node[entity] (Usuario)   at (0,   0)  {Usuario};
\node[entity] (Vendedor)  at (-6,  0)  {Vendedor};
\node[entity] (Pedido)    at (6,   0)  {Pedido};
\node[entity] (Resena)    at (0,  -5)  {Reseña};
\node[entity] (Producto)  at (-6, -5)  {Producto};
\node[entity] (Categoria) at (-6, -11)  {Categoría};

% ==========================================
% 2. RELACIONES (DIAMANTES)
% ==========================================

\node[relationship] (RelUserSeller) at ($(Usuario)!0.5!(Vendedor)$)   {Tiene};
\node[relationship] (RelUserOrder)  at ($(Usuario)!0.5!(Pedido)$)     {Realiza};
\node[relationship] (RelUserReview) at ($(Usuario)!0.5!(Resena)$)     {Escribe};
\node[relationship] (RelSellerProd) at ($(Vendedor)!0.5!(Producto)$)  {Publica};
\node[relationship] (RelOrderProd)  at ($(Pedido)!0.5!(Producto)$)    {Contiene};
\node[relationship] (RelReviewProd) at ($(Resena)!0.5!(Producto)$)    {Sobre};
\node[relationship] (RelProdCat)    at ($(Producto)!0.5!(Categoria)$) {Pertenece};

% ==========================================
% 3. CONEXIONES Y CARDINALIDAD
% ==========================================

\draw[link] (Usuario)  -- node[card, pos=0.2]{1} (RelUserSeller)
                      -- node[card, pos=0.8]{1} (Vendedor);

\draw[link] (Usuario)  -- node[card, pos=0.2]{1} (RelUserOrder)
                      -- node[card, pos=0.8]{N} (Pedido);

\draw[link] (Usuario)  -- node[card, pos=0.2]{1} (RelUserReview)
                      -- node[card, pos=0.8]{N} (Resena);

\draw[link] (Vendedor) -- node[card, pos=0.2]{1} (RelSellerProd)
                      -- node[card, pos=0.8]{N} (Producto);

\draw[link] (Pedido)   -- node[card, pos=0.2]{N} (RelOrderProd)
                      -- node[card, pos=0.8]{M} (Producto);

\draw[link] (Resena)   -- node[card, pos=0.2]{N} (RelReviewProd)
                      -- node[card, pos=0.8]{1} (Producto);

\draw[link] (Producto) -- node[card, pos=0.2]{N} (RelProdCat)
                      -- node[card, pos=0.8]{1} (Categoria);

% ==========================================
% 4. ATRIBUTOS
% ==========================================

% --- USUARIO ---
\draw[link] (Usuario) -- ++(0.7,1.5)
    node[key_attribute] {\underline{ID}};
\draw[link] (Usuario) -- ++(-3,2)
    node[attribute] {Nombre};
\draw[link] (Usuario) -- ++(2.5,1.2)
    node[attribute] {Email};
\draw[link] (Usuario) -- ++(3,2.5)
    node[attribute] {Rol};
\draw[link] (Usuario) -- ++(-1,2)
    node[attribute] {Password};

% --- VENDEDOR ---
\draw[link] (Vendedor) -- ++(1,1.5)
    node[key_attribute] {\underline{ID}};
\draw[link] (Vendedor) -- ++(-2.5,1.2)
    node[attribute] {Nombre\\Tienda};
\draw[link] (Vendedor) -- ++(-3,0)
    node[attribute] {Rating};
\draw[link] (Vendedor) -- ++(-2.5,-1.2)
    node[attribute] {Ubicación};
\draw[link] (Vendedor) -- ++(2.5,1.2)
    node[attribute] {Estado};

% --- PEDIDO ---
\draw[link] (Pedido) -- ++(1,1.5)
    node[key_attribute] {\underline{ID}};
\draw[link] (Pedido) -- ++(2.8,1.0)
    node[attribute] {Total};
\draw[link] (Pedido) -- ++(3,0)
    node[attribute] {Estado};
\draw[link] (Pedido) -- ++(2.8,-1.0)
    node[attribute] {Fecha};

% --- RESEÑA ---
\draw[link] (Resena) -- ++(1,-1.5)
    node[key_attribute] {\underline{ID}};
\draw[link] (Resena) -- ++(2.8,-1.0)
    node[attribute] {Comentario};
\draw[link] (Resena) -- ++(-2.8,-1.0)
    node[attribute] {Calificación};

% --- PRODUCTO ---
\draw[link] (Producto) -- ++(-3,0)
    node[key_attribute] {\underline{ID}};
\draw[link] (Producto) -- ++(-2.8,1.0)
    node[attribute] {Título};
\draw[link] (Producto) -- ++(-2.8,-1.0)
    node[attribute] {Precio};
\draw[link] (Producto) -- ++(-1.5,-1.5)
    node[attribute] {Stock};
\draw[link] (Producto) -- ++(-1.5,1.5)
    node[attribute] {Marca};

% --- CATEGORÍA ---
\draw[link] (Categoria) -- ++(-3,0)
    node[key_attribute] {\underline{ID}};
\draw[link] (Categoria) -- ++(3,0)
    node[attribute] {Nombre};
\draw[link] (Categoria) -- ++(0,-1.5)
    node[attribute] {Icono};

\end{tikzpicture}

\caption{Diagrama Entidad-Relación - Silicon-Trail}
\end{figure}

\end{document}